% LaTeX Template für Abgaben an der Universität Stuttgart
% Autor: Sandro Speth
% Bei Fragen: Sandro.Speth@iste.uni-stuttgart.de
% Contributor: JSendra - extended and modified for current Abgaben schema
%-----------------------------------------------------------
% Modul beinhaltet Befehl fuer Aufgabennummerierung,
% sowie die Header Informationen.

% Überschreibt enumerate Befehl, sodass 1. Ebene Items mit
\renewcommand{\theenumi}{(\alph{enumi})}
% (a), (b), etc. nummeriert werden.
\renewcommand{\labelenumi}{\text{\theenumi}}

% Counter für das Blatt und die Aufgabennummer.
% Ersetze die Nummer des Übungsblattes und die Nummer der Aufgabe
% den Anforderungen entsprechend.
% Gesetz werden die counter in der hauptdatei, damit siese hier nicht jedes mal verändert werden muss
% Beachte:
% \setcounter{countername}{number}: Legt den Wert des Counters fest
% \stepcounter{countername}: Erhöht den Wert des Counters um 1.
\newcounter{sheetnr}
\newcounter{exnum}

\newcounter{teilaufgabe}[exnum]


% Befehl für die Aufgabentitel
\newcommand{\exercise}[1]{
    \section*{
        Aufgabe \theexnum\stepcounter{exnum}: #1
        }
} % Befehl für Aufgabentitel

\newcommand{\teil}[1]{
    \section*{
        Aufgabe \theexnum.\theteilaufgabe\stepcounter{teilaufgabe}: #1
    }
}


% Formatierung der Kopfzeile
% \ohead: Setzt rechten Teil der Kopfzeile mit
% Namen und Matrikelnummern aller Bearbeiter
\ohead{Josse Sendra (3542251)\\
Dennis Hehn (3593682)\\
Matthias Wagner (3542811)}
% \chead{} kann mittleren Kopfzeilen Teil sezten
% \ihead: Setzt linken Teil der Kopfzeile mit
% Modulnamen, Semester und Übungsblattnummer
\ihead{NumStatStoch\\
Sommersemester 26\\
Blatt \thesheetnr}


